% Options for packages loaded elsewhere
% Options for packages loaded elsewhere
\PassOptionsToPackage{unicode}{hyperref}
\PassOptionsToPackage{hyphens}{url}
%
\documentclass[
]{book}
\usepackage{xcolor}
\usepackage{amsmath,amssymb}
\setcounter{secnumdepth}{-\maxdimen} % remove section numbering
\usepackage{iftex}
\ifPDFTeX
  \usepackage[T1]{fontenc}
  \usepackage[utf8]{inputenc}
  \usepackage{textcomp} % provide euro and other symbols
\else % if luatex or xetex
  \usepackage{unicode-math} % this also loads fontspec
  \defaultfontfeatures{Scale=MatchLowercase}
  \defaultfontfeatures[\rmfamily]{Ligatures=TeX,Scale=1}
\fi
\usepackage{lmodern}
\ifPDFTeX\else
  % xetex/luatex font selection
\fi
% Use upquote if available, for straight quotes in verbatim environments
\IfFileExists{upquote.sty}{\usepackage{upquote}}{}
\IfFileExists{microtype.sty}{% use microtype if available
  \usepackage[]{microtype}
  \UseMicrotypeSet[protrusion]{basicmath} % disable protrusion for tt fonts
}{}
\makeatletter
\@ifundefined{KOMAClassName}{% if non-KOMA class
  \IfFileExists{parskip.sty}{%
    \usepackage{parskip}
  }{% else
    \setlength{\parindent}{0pt}
    \setlength{\parskip}{6pt plus 2pt minus 1pt}}
}{% if KOMA class
  \KOMAoptions{parskip=half}}
\makeatother
% Make \paragraph and \subparagraph free-standing
\makeatletter
\ifx\paragraph\undefined\else
  \let\oldparagraph\paragraph
  \renewcommand{\paragraph}{
    \@ifstar
      \xxxParagraphStar
      \xxxParagraphNoStar
  }
  \newcommand{\xxxParagraphStar}[1]{\oldparagraph*{#1}\mbox{}}
  \newcommand{\xxxParagraphNoStar}[1]{\oldparagraph{#1}\mbox{}}
\fi
\ifx\subparagraph\undefined\else
  \let\oldsubparagraph\subparagraph
  \renewcommand{\subparagraph}{
    \@ifstar
      \xxxSubParagraphStar
      \xxxSubParagraphNoStar
  }
  \newcommand{\xxxSubParagraphStar}[1]{\oldsubparagraph*{#1}\mbox{}}
  \newcommand{\xxxSubParagraphNoStar}[1]{\oldsubparagraph{#1}\mbox{}}
\fi
\makeatother


\usepackage{longtable,booktabs,array}
\usepackage{calc} % for calculating minipage widths
% Correct order of tables after \paragraph or \subparagraph
\usepackage{etoolbox}
\makeatletter
\patchcmd\longtable{\par}{\if@noskipsec\mbox{}\fi\par}{}{}
\makeatother
% Allow footnotes in longtable head/foot
\IfFileExists{footnotehyper.sty}{\usepackage{footnotehyper}}{\usepackage{footnote}}
\makesavenoteenv{longtable}
\usepackage{graphicx}
\makeatletter
\newsavebox\pandoc@box
\newcommand*\pandocbounded[1]{% scales image to fit in text height/width
  \sbox\pandoc@box{#1}%
  \Gscale@div\@tempa{\textheight}{\dimexpr\ht\pandoc@box+\dp\pandoc@box\relax}%
  \Gscale@div\@tempb{\linewidth}{\wd\pandoc@box}%
  \ifdim\@tempb\p@<\@tempa\p@\let\@tempa\@tempb\fi% select the smaller of both
  \ifdim\@tempa\p@<\p@\scalebox{\@tempa}{\usebox\pandoc@box}%
  \else\usebox{\pandoc@box}%
  \fi%
}
% Set default figure placement to htbp
\def\fps@figure{htbp}
\makeatother





\setlength{\emergencystretch}{3em} % prevent overfull lines

\providecommand{\tightlist}{%
  \setlength{\itemsep}{0pt}\setlength{\parskip}{0pt}}



 


\usepackage[utf8]{inputenc}
\usepackage{dnd}
\usepackage{dndbook}
\makeatletter
\@ifpackageloaded{caption}{}{\usepackage{caption}}
\AtBeginDocument{%
\ifdefined\contentsname
  \renewcommand*\contentsname{Table of contents}
\else
  \newcommand\contentsname{Table of contents}
\fi
\ifdefined\listfigurename
  \renewcommand*\listfigurename{List of Figures}
\else
  \newcommand\listfigurename{List of Figures}
\fi
\ifdefined\listtablename
  \renewcommand*\listtablename{List of Tables}
\else
  \newcommand\listtablename{List of Tables}
\fi
\ifdefined\figurename
  \renewcommand*\figurename{Figure}
\else
  \newcommand\figurename{Figure}
\fi
\ifdefined\tablename
  \renewcommand*\tablename{Table}
\else
  \newcommand\tablename{Table}
\fi
}
\@ifpackageloaded{float}{}{\usepackage{float}}
\floatstyle{ruled}
\@ifundefined{c@chapter}{\newfloat{codelisting}{h}{lop}}{\newfloat{codelisting}{h}{lop}[chapter]}
\floatname{codelisting}{Listing}
\newcommand*\listoflistings{\listof{codelisting}{List of Listings}}
\makeatother
\makeatletter
\makeatother
\makeatletter
\@ifpackageloaded{caption}{}{\usepackage{caption}}
\@ifpackageloaded{subcaption}{}{\usepackage{subcaption}}
\makeatother
\usepackage{bookmark}
\IfFileExists{xurl.sty}{\usepackage{xurl}}{} % add URL line breaks if available
\urlstyle{same}
\hypersetup{
  pdftitle={The Shadow of Mars},
  pdfauthor={Campaign Module},
  hidelinks,
  pdfcreator={LaTeX via pandoc}}


\title{The Shadow of Mars}
\usepackage{etoolbox}
\makeatletter
\providecommand{\subtitle}[1]{% add subtitle to \maketitle
  \apptocmd{\@title}{\par {\large #1 \par}}{}{}
}
\makeatother
\subtitle{A Roman Campaign in Five Acts}
\author{Campaign Module}
\date{2026-01-18}
\begin{document}
\frontmatter
\maketitle


\mainmatter
\chapter{Introduction}\label{introduction}

\section{Campaign Overview}\label{campaign-overview}

\emph{The Shadow of Mars} is a 5-session campaign set in a fantastical
Roman world where the gods are real, auguries matter, and destiny can be
changed---or embraced.

The party plays as members of a military unit stationed at a remote
frontier fort. When dark omens appear and a powerful artifact is
unearthed, they must decide whether to serve Rome, the gods, or their
own ambitions.

\textbf{Levels:} Characters start at level 3 and advance to level 7 by
the campaign's end.

\textbf{Themes:} Duty vs.~ambition, fate vs.~free will, loyalty
vs.~survival, divine intervention

\textbf{Setting:} The campaign takes place along Rome's Germanic
frontier in 175 AD, in an alternate world where D\&D magic exists
alongside Roman military discipline and pantheon worship.

\section{Key NPCs}\label{key-npcs}

\begin{itemize}
\tightlist
\item
  \textbf{Legate Marcus Aurelius Corvinus} -- Commander of the fort,
  ambitious and ruthless
\item
  \textbf{Augur Cassia Liviana} -- The fort's seer, troubled by dark
  visions
\item
  \textbf{Centurion Titus Varro} -- Veteran soldier, loyal but pragmatic
\item
  \textbf{Tribune Lucius Valerius Maximus} -- A mysterious figure from
  Rome with a hidden agenda
\item
  \textbf{Vercingetorix the Red} -- Germanic chieftain and reluctant
  ally
\item
  \textbf{Senator Gaius Cassius Brutus} -- Roman aristocrat with
  revolutionary ambitions
\end{itemize}

\section{NPC OGAS Framework (Guy Sclanders
Method)}\label{npc-ogas-framework-guy-sclanders-method}

\subsection{Legate Marcus Aurelius
Corvinus}\label{legate-marcus-aurelius-corvinus}

\textbf{Objective:} Secure a military triumph and consulship in Rome
\textbf{Goal:} Possess the Spear of Mars as proof of his frontier
victory \textbf{Agenda:} Use the artifact to gain Senate favor and
advancement \textbf{Secret:} He's drowning in gambling debts to a
powerful Senator (Brutus) who now owns him

\textbf{Personality:} Blunt, aggressive, views everything through
military lens \textbf{Speaks like:} A drill sergeant - short, commanding
sentences \textbf{Reacts to opposition:} With anger first, calculation
second \textbf{If threatened:} Will sacrifice subordinates to save
himself

\subsection{Augur Cassia Liviana}\label{augur-cassia-liviana}

\textbf{Objective:} Prevent divine catastrophe and protect Rome from
Mars' wrath \textbf{Goal:} Destroy the Spear of Mars before it corrupts
anyone else \textbf{Agenda:} Guide the party toward making the right
choice without forcing them \textbf{Secret:} She's Mars' daughter
(demigod) and can feel the spear calling to her bloodline

\textbf{Personality:} Quiet, intense, speaks in riddles when nervous
\textbf{Speaks like:} An oracle - cryptic, layered meanings, pauses
between thoughts \textbf{Reacts to opposition:} With sadness rather than
anger; tries to show visions \textbf{If threatened:} Will sacrifice
herself to stop the spear

\subsection{Centurion Titus Varro}\label{centurion-titus-varro}

\textbf{Objective:} Survive his final year of service and retire with
honor \textbf{Goal:} Keep his soldiers alive and maintain order at the
fort \textbf{Agenda:} Support whoever seems most competent---party or
commander \textbf{Secret:} He knows the Legate is corrupt but can't
prove it; 20 years of service make him complicit

\textbf{Personality:} Pragmatic, weary, dark humor to cope with violence
\textbf{Speaks like:} A veteran NCO - profane, direct, occasional wisdom
\textbf{Reacts to opposition:} With exhausted reason; ``I've seen this
before'' \textbf{If threatened:} Will choose his men over his commander
every time

\subsection{Tribune Lucius Valerius
Maximus}\label{tribune-lucius-valerius-maximus}

\textbf{Objective:} Earn respect as a ``real Roman'' despite his
aristocratic birth \textbf{Goal:} Complete his secret mission from
Senator Brutus (destroy the spear) \textbf{Agenda:} Manipulate the party
into destroying it so he gets credit \textbf{Secret:} He's Senator
Brutus' illegitimate son, desperate to prove his worth

\textbf{Personality:} Overconfident, theatrical, masks insecurity with
arrogance \textbf{Speaks like:} An aristocrat - flowery language, quotes
philosophers, condescending \textbf{Reacts to opposition:} With wounded
pride; becomes petty and vindictive \textbf{If threatened:} Will betray
anyone to survive; coward beneath the bravado

\subsection{Vercingetorix the Red}\label{vercingetorix-the-red}

\textbf{Objective:} Protect his people from Roman expansion
\textbf{Goal:} Recover the Spear of Mars (killed his grandfather,
tribe's sacred enemy) \textbf{Agenda:} Use the party to get past Roman
defenses, then take the spear by force or alliance \textbf{Secret:} He's
dying from a wasting disease; this is his last campaign; wants his death
to mean something

\textbf{Personality:} Noble warrior, tired of war, respects strength and
honesty \textbf{Speaks like:} Poetic warrior - comparisons to nature,
warrior's code \textbf{Reacts to opposition:} Tests strength first,
respects those who stand firm \textbf{If threatened:} Will fight to the
death but prefers honorable single combat

\subsection{Senator Gaius Cassius
Brutus}\label{senator-gaius-cassius-brutus}

\textbf{Objective:} ``Restore the Republic'' (overthrow the Emperor,
install himself as First Citizen) \textbf{Goal:} Acquire the Spear of
Mars to assassinate the Emperor during his Triumph \textbf{Agenda:}
Manipulate everyone as pawns; eliminate witnesses once spear is secured
\textbf{Secret:} He's already made pacts with three other gods (Jupiter,
Pluto, Bellona) promising them Mars' power if they support his coup

\textbf{Personality:} Charismatic sociopath, genuinely believes he's
Rome's savior \textbf{Speaks like:} Statesman - rhetorical questions,
appeals to Roman values, very persuasive \textbf{Reacts to opposition:}
With charm first, then legal/political pressure, then assassination
\textbf{If threatened:} Flees and regroups; will burn Rome before
admitting failure

\subsection{Mars (God of War)}\label{mars-god-of-war}

\textbf{Objective:} Reclaim his honor after mortals stole and destroyed
his divine weapon \textbf{Goal:} Make the party understand the
consequences of defying a god \textbf{Agenda:} Test their worth---if
strong, make them his champions; if weak, destroy them \textbf{Secret:}
He's secretly proud they destroyed the corrupted spear; it had become a
blight on his name

\textbf{Personality:} Wrathful but honorable, respects courage above all
\textbf{Speaks like:} God of war - imperious, archaic Latin mixed with
English, SHORT and LOUD \textbf{Reacts to opposition:} With fury but
also respect; violence is prayer to him \textbf{If threatened:} Cannot
truly die; will curse them across lifetimes if dishonored

\section{Using OGAS in Play}\label{using-ogas-in-play}

\textbf{When NPCs appear:} 1. Start with their \textbf{Goal} - what they
want THIS SESSION 2. Use their \textbf{Agenda} - HOW they'll try to get
it 3. Reveal \textbf{Objective} through multiple sessions 4. Drop
\textbf{Secret} hints for clever players to discover

\textbf{Example: When party meets Legate (Session 1)} - Goal: Get them
to retrieve spear → Direct orders - Agenda: Establish authority →
Threatens court-martial if they refuse - Objective leak: ``Securing this
artifact will echo all the way to Rome'' - Secret hint: Cassia whispers,
``The Legate is not his own master''

\textbf{If party goes off-script:} NPCs pursue their OGAS regardless of
party actions. This makes world feel reactive and alive.

\section{The Central Mystery}\label{the-central-mystery}

A cursed artifact called the \emph{Spear of Mars} has been unearthed
near the fort. Legend says it once belonged to the war god himself and
was used to kill a divine being. Now it whispers to those who touch it,
promising power and glory while slowly corrupting their souls.

The party must decide: destroy it, claim it, or deliver it to Rome?

\chapter{Session 0 Questions}\label{session-0-questions}

\chapter{Session 0 -- D\&D 5e (5-Session
Module)}\label{session-0-dd-5e-5-session-module}

\section{Campaign Basics}\label{campaign-basics}

\begin{itemize}
\tightlist
\item
  Are you okay with this being a short campaign (about 5 sessions) with
  a clear ending?
\item
  Are character death and permanent consequences on the table?
\item
  Do you prefer more combat, more roleplay, or a balance of both?
\item
  Are you okay with the world reacting strongly to your choices?
\end{itemize}

\section{Tone \& Comfort}\label{tone-comfort}

\begin{itemize}
\tightlist
\item
  What tone do you want: serious, gritty, heroic, or light?
\item
  Are there any themes or topics you want to avoid?
\item
  Is PvP (stealing from or attacking party members) allowed?
\item
  Are you okay with NPC authority figures giving orders?
\end{itemize}

\section{Characters \& Party}\label{characters-party}

\begin{itemize}
\tightlist
\item
  Do your characters already know each other, or do they meet in session
  1?
\item
  Why does your character choose to work with this group?
\item
  What is one clear strength your character brings to the party?
\item
  What is one weakness or flaw your character has?
\end{itemize}

\section{World Assumptions}\label{world-assumptions}

\begin{itemize}
\tightlist
\item
  Are you okay with law, rank, and social status mattering in the world?
\item
  Are you okay with consequences for breaking the law?
\item
  Should reputation follow your characters?
\end{itemize}

\section{Combat \& Play Style}\label{combat-play-style}

\begin{itemize}
\tightlist
\item
  Do you prefer combat to be fast and cinematic or slower and tactical?
\item
  Are you okay with enemies retreating or surrendering?
\item
  Is retreat a valid option for the party?
\item
  Are you okay with teamwork and positioning mattering in combat?
\end{itemize}

\section{Magic \& Gods}\label{magic-gods}

\begin{itemize}
\tightlist
\item
  Should magic feel common, rare, or feared?
\item
  Should the gods be clearly real or distant and mysterious?
\item
  Are omens and divine signs something you want in the game?
\end{itemize}

\section{Table Rules}\label{table-rules}

\begin{itemize}
\tightlist
\item
  Do we play rules as written, rules as intended, or rule of cool?
\item
  How should we handle rule disputes during play?
\item
  Are there any table rules or expectations you want to add?
\end{itemize}

\section{End of Module}\label{end-of-module}

\begin{itemize}
\tightlist
\item
  Are you okay if your character's story ends in death, promotion,
  disgrace, or major change?
\item
  What kind of ending do you hope for: victorious, bittersweet, tragic,
  or uncertain?
\end{itemize}

\section{Final Check}\label{final-check}

\begin{itemize}
\tightlist
\item
  Is there anything you need from me as DM to enjoy this campaign?
\item
  Is there anything you are concerned about before we start?
\item
  Is everyone committed to playing the full module?
\end{itemize}

\chapter{Chapter 1: Blood and Omens}\label{chapter-1-blood-and-omens}

\textbf{Applying Guy Sclanders' Methodology} \emph{(Based on
\href{https://books.google.com/books/about/The_Complete_Guide_to_Creating_Epic_Camp.html?id=0XFNvgEACAAJ}{The
Complete Guide to Creating Epic Campaigns} and
\href{https://www.greatgamemaster.com/dm/about-us/}{How to be a Great
Game Master})}

\section{Session Overview}\label{session-overview}

The party arrives at Fort Vindolanda on the Germanic frontier. Strange
omens plague the camp, and the party is sent to investigate a nearby
excavation where something ancient has been unearthed.

\textbf{Level:} 3 → Advance to 4 at session end \textbf{Duration:} 3-4
hours \textbf{Key Goal:} Investigate the excavation site and decide the
spear's fate \textbf{Theme:} Duty vs.~Free Will (first exploration)

\subsection{Guy Sclanders' Principles
Applied}\label{guy-sclanders-principles-applied}

\textbf{Pressure Cooker Design:} - Time pressure: The spear is
corrupting the fort NOW - Incomplete information: Who can be trusted? -
Escalating stakes: Personal → Fort → Rome → Divine

\textbf{OGAS Integration:} Legate, Augur, and Centurion all have clear
goals THIS SESSION

\textbf{Meaningful Choices:} No ``correct'' answer; every choice has
consequences

\section{Pre-Session Preparation}\label{pre-session-preparation}

\subsection{Props \& Handouts (Physical
Immersion)}\label{props-handouts-physical-immersion}

\textbf{Essential Props:} 1. \textbf{Wax Tablet} - Write the Legate's
orders in Latin on a small board 2. \textbf{Red Dice} - For corruption
tracking (give corrupted players red d6s) 3. \textbf{Omen Cards} - Index
cards with Latin phrases describing bad omens 4. \textbf{Spear Token} -
Black painted dowel to represent the spear physically

\textbf{Handout 1: Legate's Orders (Latin/English)}

\begin{verbatim}
MANDATUM LEGATI MARCI AURELII CORVINI

Ad milites electos:
Descendite in ruinas. Recuperate telum divinum.
Nemo tangat telum nisi ego permitto.

(To the chosen soldiers:
Descend into the ruins. Recover the divine weapon.
No one shall touch the weapon unless I permit it.)

In nomine Senatus Populique Romani.
\end{verbatim}

\textbf{Handout 2: Gaulish Runes Translation}

\begin{verbatim}
GAULISH/CELTIC INSCRIPTION:
"Atir nertos, atir bitus, atir maros"
(Father of strength, father of life, father of death)

"Ro·biđ·su in gaissos·mo dîvon-macon"
(Here lies the spear, god-slayer)

TRANSLATION:
Here lies the God-Killer, the Spear of Broken Oaths,
sealed by the blood of heroes. May it never taste sunlight again.
Forged in vengeance. Quenched in divine blood.
The son of Mars fell here. His father's wrath lingers still.
\end{verbatim}

\textbf{Handout 3: Omens Observed (Give to players as prop)} - All camp
dogs died at midnight (no visible cause) - Sacrificial sheep had black
livers (sign of divine curse) - Ravens fly backwards over the principia
- Springs run red with iron-tainted water - Sentries report shadows
moving without bodies - Legionaries hear distant war-horns (no enemy
visible)

\section{Scene 0: Cold Open (5
minutes)}\label{scene-0-cold-open-5-minutes}

\textbf{Purpose:} Hook players immediately with visceral horror (Guy
Sclanders: ``Strong Opening'')

\textbf{Start in media res - No preamble:}

\textbf{Read aloud (speak slowly, with dread):} \textgreater{}
\emph{Three nights ago.} \textgreater{} \textgreater{} You dream. All of
you, the same dream. You stand in darkness, ankle-deep in blood that
doesn't ripple. In the distance, a spear stands upright in the earth,
glowing like a dying star. When you try to look away, you can't. When
you try to wake, you can't. \textgreater{} \textgreater{} A voice---old,
vast, furious---speaks a single word in a language you shouldn't
understand, but do: \textgreater{} \textgreater{} \textbf{``MEUM.''
(MINE.)} \textgreater{} \textgreater{} You wake. Your tent reeks of
copper and smoke. Outside, the camp dogs are howling. \textgreater{}
\textgreater{} By dawn, every dog in Fort Vindolanda is dead.

\textbf{DM Direction:} Don't answer questions yet. Cut to black, then:

\textbf{Read aloud:} \textgreater{} \emph{Now.} \textgreater{}
\textgreater{} The principia. The Legate's headquarters. The air is
thick with incense, but underneath it, you smell rot. Around the table:
Legate Corvinus, eyes bloodshot and intense. Augur Cassia, pale as
death, clutching a bloodstained cloth. Centurion Varro, one hand on his
gladius, looking like he hasn't slept. \textgreater{} \textgreater{} The
Legate speaks. His voice is iron scraping stone: \textgreater{}
\textgreater{} ``You've been chosen. Not for honor. For survival.
Something is beneath this fort. Something that killed sixty dogs in one
night. Something that made the augurs vomit blood during morning
sacrifice. We've lost four men already---two to madness, two to each
other's blades.'' \textgreater{} \textgreater{} He unrolls a rough
sketch of underground ruins. \textgreater{} \textgreater{} ``You're
going down there. You're bringing it up. You're not asking questions.''
\textgreater{} \textgreater{} Across the table, the Augur whispers:
``\emph{Mars aversus est.} The god of war has turned his face from us.''

\textbf{Pressure Cooker Element:} Players don't control the situation
yet---they're already deep in crisis

\section{Scene 1: Meeting the Legate}\label{scene-1-meeting-the-legate}

The party is summoned to meet \textbf{Legate Marcus Aurelius Corvinus}
in the principia. He's a hard man in his forties, scarred and ambitious.

\textbf{Legate's Briefing:} - Three days ago, a work crew digging a new
well discovered ancient ruins beneath the fort - They unearthed an iron
spear covered in strange runes - That night, all the camp dogs died -
The next day, the augurs found corrupted omens: black livers, reversed
bird flights - One worker touched the spear and went mad, killing two
companions before being subdued - The Legate wants the party to go down
into the ruins, secure the artifact, and bring it to him

\textbf{Complications:} - \textbf{Augur Cassia Liviana} interrupts,
warning that the omens forbid disturbing the spear - \textbf{Centurion
Varro} suggests destroying it immediately - The Legate overrules them
both---Rome must possess this weapon

\textbf{Decision Point:} Will the party obey orders or heed the augur's
warning?

\section{Scene 2: Descent into the
Ruins}\label{scene-2-descent-into-the-ruins}

The excavation site is a crude shaft dug through Roman stonework into
older Germanic structures. The air grows cold as the party descends.

\textbf{Exploration:} - Ancient Germanic temple dedicated to a forgotten
war god - Walls covered in bronze age runes (DC 15 History or Religion
to read) - Runes tell a story: \emph{``Here lies the God-Killer, the
Spear of Broken Oaths, sealed by the blood of heroes. May it never taste
sunlight again.''} - Three chambers: 1. \textbf{Hall of Shields} --
Dozens of ancient shields line the walls, some bearing symbols of
forgotten tribes 2. \textbf{Chamber of Chains} -- Iron chains hang from
ceiling, sized for something enormous 3. \textbf{The Vault} -- Central
chamber where the spear was found, now guarded by \textbf{animated
armor} (2x) possessed by vengeful spirits

\textbf{Combat Encounter:} 2 Animated Armor (CR 1 each) - The spirits
scream: ``DO NOT TAKE THE SPEAR!'' - If defeated: They dissolve,
whispering ``You have doomed us all''

\section{Scene 3: The Spear}\label{scene-3-the-spear}

The \textbf{Spear of Mars} rests on a stone altar, surrounded by broken
chains. It's eight feet long, made of black iron, covered in glowing red
runes.

\textbf{Mechanics:} - Anyone who touches it must make a DC 14 Wisdom
saving throw - \textbf{Failure:} Visions of glory and conquest flood
their mind. They gain 1 level of corruption (track this) -
\textbf{Success:} They resist but feel its pull

\textbf{Properties (if attuned):} - +2 magic weapon - Deals an extra 1d8
necrotic damage - Disadvantage on Wisdom saving throws - User becomes
increasingly aggressive and paranoid

\textbf{Roleplaying:} The spear whispers promises: - ``You could be a
hero'' - ``Rome will kneel before you'' - ``Your enemies will break like
wheat''

\section{Scene 4: Return to the
Surface}\label{scene-4-return-to-the-surface}

As the party emerges, they find the fort in chaos: - The mad worker has
escaped his bonds - He's killed three guards and is rampaging through
the camp - He's screaming: ``THE SPEAR CALLS FOR BLOOD!''

\textbf{Combat Encounter:} Corrupted Worker (use \textbf{Berserker} stat
block, CR 2) - If the party has the spear, he focuses on whoever carries
it - Dies screaming: ``Mars\ldots demands\ldots sacrifice\ldots{}''

\section{Conclusion}\label{conclusion}

The Legate demands the spear. The Augur begs them to destroy it.
Centurion Varro looks to the party for guidance.

\textbf{Key Questions:} - Do they give the spear to the Legate? - Do
they try to destroy it? - Does anyone attune to it?

\textbf{Advance to Level 4}

\section{DM Notes}\label{dm-notes}

\subsection{Roman Military Ranks (for
reference)}\label{roman-military-ranks-for-reference}

\begin{itemize}
\tightlist
\item
  \textbf{Legate} -- Commander of the legion (like a general)
\item
  \textbf{Tribune} -- Staff officer from Rome (aristocrat)
\item
  \textbf{Centurion} -- Company commander (60-80 men)
\item
  \textbf{Optio} -- Second-in-command to centurion
\item
  \textbf{Legionary} -- Standard soldier
\end{itemize}

\subsection{Corruption Mechanic}\label{corruption-mechanic}

Track corruption levels (0-5) for any character who touches or attunes
to the spear: - \textbf{Level 1:} Disturbing dreams of conquest -
\textbf{Level 2:} Disadvantage on Wisdom saves - \textbf{Level 3:} Must
succeed DC 12 Wisdom save to retreat from combat - \textbf{Level 4:}
Must succeed DC 14 Wisdom save to avoid attacking allies when bloodied -
\textbf{Level 5:} Fully corrupted, becomes NPC under DM control

\subsection{Optional Encounters}\label{optional-encounters}

\begin{itemize}
\tightlist
\item
  \textbf{Scouts' Report:} Germanic tribe activity increased in the
  forest
\item
  \textbf{Camp Tensions:} Legionaries are terrified, some talk of
  desertion
\item
  \textbf{The Augur's Warning:} Private conversation revealing that the
  spear belongs to Mars and was used in a divine war \# Chapter 2: The
  Tribune's Gambit
\end{itemize}

\section{Session Overview}\label{session-overview-1}

A Tribune arrives from Rome with secret orders regarding the spear.
Meanwhile, Germanic raiders attack the fort, and the party must defend
it while uncovering the Tribune's true agenda.

\textbf{Level:} 4 → Advance to 5 at session end \textbf{Key Goal:}
Defend the fort and discover the Tribune's mission

\section{Opening Scene: The Tribune
Arrives}\label{opening-scene-the-tribune-arrives}

Three days after the spear's discovery, a Tribune arrives from Rome with
an elite escort of Praetorian Guards. He's young, polished, and
dangerously intelligent.

\textbf{Read aloud:} \textgreater{} The gates open to admit a small
column of riders in gleaming armor. At their head rides a man too clean
for the frontier---purple-trimmed toga, oiled hair, rings on every
finger. He surveys the fort like a merchant appraising damaged goods.
Behind him, six Praetorians sit their horses like statues of iron and
menace.

\textbf{Tribune Lucius Valerius Maximus} introduces himself with
imperial authority: - He carries sealed orders from the Senate - The
spear is to be delivered to Rome immediately - The Legate is relieved of
command pending investigation - The party is promoted to ``Special
Envoys of the Senate'' (higher authority than the Legate) - They will
escort the Tribune and the spear to Rome

\textbf{Complications:} - The Legate is furious but powerless - The
Augur warns this is exactly what Mars wants---the spear in Rome -
Centurion Varro quietly tells the party: ``That Tribune is lying about
something''

\section{Scene 1: Investigation \&
Tensions}\label{scene-1-investigation-tensions}

The party has time to investigate before departure. Options:

\textbf{Talk to the Augur (Cassia Liviana):} - She's had visions: the
spear in Rome, the Capitoline Hill running with blood - She reveals the
spear's true history: it was used by a mortal hero to kill a minor god
(Mars' son) - Mars cursed it so that whoever wields it will be consumed
by war - The Tribune knows this---he plans to use it to assassinate
someone in Rome

\textbf{Talk to Centurion Varro:} - He's served 20 years and knows
political games - ``That Tribune didn't come here for the spear---he
came here to disappear the spear'' - He suspects the Tribune plans to
fake an attack and blame the party - Offers to help if things go wrong

\textbf{Search the Tribune's Quarters (DC 16 Investigation):} - Find a
letter (encoded, DC 14 Intelligence to decipher) - The letter is from a
Senator ordering the Tribune to ``ensure the artifact never reaches
Rome'' - The Tribune is supposed to destroy it---or those carrying it

\section{Scene 2: The Raid}\label{scene-2-the-raid}

That night, Germanic raiders attack the fort in force. This is
coordinated with the Tribune's arrival (suspicious timing).

\textbf{Battle Map:} Fort under siege - 20+ raiders (use \textbf{Tribal
Warrior} stat blocks in waves) - 2 \textbf{Berserkers} (CR 2 each) - 1
\textbf{Germanic Chieftain} - \textbf{Vercingetorix the Red} (use
\textbf{Gladiator} stat block, CR 5)

\textbf{Combat Phases:}

\textbf{Phase 1: The Walls (2 rounds)} - Raiders swarm the palisade with
ladders - Party can man ballistae (2d10 piercing, DC 12 Dex to aim) -
Archers provide supporting fire

\textbf{Phase 2: Breach (3 rounds)} - Raiders break through the gate
with a battering ram - Melee combat in the fort courtyard - Centurion
Varro fights alongside the party

\textbf{Phase 3: The Headquarters (2 rounds)} - Vercingetorix and
berserkers push toward where the spear is kept - The Tribune and his
Praetorians mysteriously don't help - Climactic battle to defend the
principia

\section{Scene 3: The Truth Revealed}\label{scene-3-the-truth-revealed}

After the battle, as the party catches their breath:

\textbf{The Tribune's Move:} - He accuses the party of being traitors
who let the raiders in - Claims they were trying to steal the spear -
Orders his Praetorians to arrest them - The Legate, humiliated and
angry, supports this

\textbf{Centurion Varro's Intervention:} - He steps forward: ``The
Tribune is the traitor!'' - Reveals what he saw: the Tribune signaling
the raiders from the walls - The Praetorians hesitate---they respect
Varro as a veteran

\textbf{Decision Point:} The party must choose: 1. \textbf{Fight} --
Combat with Tribune + 6 Praetorians (deadly) 2. \textbf{Diplomacy} --
Convince the Praetorians (DC 16 Persuasion with evidence) 3.
\textbf{Flee} -- Escape with the spear into the Germanic forest 4.
\textbf{Surrender} -- Give up the spear to see what happens

\section{Conclusion Options}\label{conclusion-options}

\textbf{If they expose the Tribune:} - He confesses under pressure: he
was ordered to ensure the spear never reached Rome - A powerful Senator
fears what it could do in the wrong hands - He offers them a deal: help
him destroy it, and they'll all walk away free - The Augur supports
destroying it

\textbf{If they flee:} - They're now fugitives from Rome with a cursed
spear - Vercingetorix intercepts them in the forest - He offers
alliance: ``My enemy's enemy is my friend''

\textbf{If they fight:} - Chaotic battle, possible deaths - They must
flee afterward - Now hunted by both Rome and the Germanic tribes

\textbf{Advance to Level 5}

\section{DM Notes}\label{dm-notes-1}

\subsection{Vercingetorix the Red (Germanic
Chieftain)}\label{vercingetorix-the-red-germanic-chieftain}

\begin{itemize}
\tightlist
\item
  \textbf{Motivation:} He's not evil---he's defending his homeland from
  Roman expansion
\item
  Can become an unlikely ally if the party treats him with respect
\item
  Has his own history with the spear: it killed his grandfather
\end{itemize}

\subsection{The Tribune's Secret}\label{the-tribunes-secret}

Lucius Valerius is caught between loyalty to Rome and orders from a
corrupt Senator. He's not wholly evil---he genuinely believes destroying
the spear is best for Rome. This creates moral complexity.

\subsection{Pacing}\label{pacing}

This session should feel tense and chaotic. The raid happens fast, and
the party must make quick decisions with incomplete
information---classic Guy Sclanders ``pressure cooker'' design. \#
Chapter 3: Through the Dark Forest

\section{Session Overview}\label{session-overview-2}

The party must journey through the Teutoburg Forest---either fleeing
Rome or on a secret mission to destroy the spear. They face natural
dangers, Germanic tribes, and the spear's corrupting influence.

\textbf{Level:} 5 → Advance to 6 at session end \textbf{Key Goal:} Reach
the sacred grove where the spear can be destroyed

\section{Opening: The Choice is Made}\label{opening-the-choice-is-made}

Based on Session 2's ending, the party now has a clear goal: -
\textbf{If allied with Tribune:} Secret mission to destroy the spear at
an ancient grove - \textbf{If fugitives:} Hunted by Rome, seeking refuge
with Vercingetorix - \textbf{If holding spear:} Fighting its corrupting
whispers every day

\textbf{Read aloud:} \textgreater{} The forest swallows you whole.
Ancient oaks rise like cathedral columns, their canopy so thick that
even at midday, twilight reigns. The Romans call this the Dark
Forest---the place where legions disappear. Your guide
(Tribune/Vercingetorix/Augur) moves carefully, watching the shadows.
Behind you, somewhere in the green gloom, pursuers gather.

\section{Scene 1: The Haunted Grove}\label{scene-1-the-haunted-grove}

The party camps in a grove that feels wrong---too quiet, too still.

\textbf{That night:} - Anyone attuned to the spear has vivid nightmares
(Wisdom save DC 15 or gain exhaustion) - Visions of ancient battles,
blood-soaked altars, screaming gods - The spear whispers: ``You could
end this. Take me to Rome. Rule the world.''

\textbf{Corruption Check:} - Anyone who touched the spear must make a
Wisdom save (DC 14) - Failure: Corruption level increases by 1 - At
corruption level 3+, they begin arguing for keeping the spear

\textbf{Morning Discovery:} - Three legionaries are found dead, throats
cut - No tracks, no signs of struggle - One has carved into a tree:
``MARS SEES ALL''

\textbf{Investigation (DC 14 Survival):} - The killer is among the
party's NPCs - One of them (your choice) is possessed by Mars' influence
through the spear - They'll try to kill whoever plans to destroy it

\section{Scene 2: The Germanic
Ambush}\label{scene-2-the-germanic-ambush}

Germanic scouts have tracked the party. Depending on their relationship
with Vercingetorix:

\textbf{If Allied with Vercingetorix:} - His scouts find them, offer
safe passage - But they demand a toll: information about Roman troop
movements - Moral dilemma: betray Rome or fight your allies?

\textbf{If Not Allied:} - Ambush by 12 \textbf{Tribal Warriors} + 2
\textbf{Scouts} (CR 1/2 each) - Combat in dense forest (difficult
terrain, limited visibility) - Can be ended via diplomacy if they drop
weapons and negotiate

\textbf{Key NPC:} \textbf{Thusnelda}, a Germanic druid - She knows of
the spear: ``The God-Killer walks again'' - Offers to guide them to the
sacred grove where it can be unmade - Warns: ``The grove demands a
sacrifice---blood for blood, life for death''

\section{Scene 3: The River Crossing}\label{scene-3-the-river-crossing}

The party must cross a swollen river to reach the sacred grove.

\textbf{Challenges:} - River is 100 feet wide, fast-flowing (DC 14
Athletics to swim) - Fallen tree serves as a bridge (DC 12 Acrobatics to
cross safely) - Rope can be strung across with teamwork

\textbf{Complication: The Hunters Arrive}

As they cross, Roman pursuers catch up: - 8 \textbf{Legionaries} (use
\textbf{Guard} stat blocks) - 1 \textbf{Centurion} (use \textbf{Knight}
stat block, CR 3) - Led by the Legate if he survived Session 2

\textbf{Set-piece Combat:} - Half the party on each side of the river -
Ranged attacks across the water - Melee on the bridge - Environmental
hazards: fast current, slippery logs

\textbf{Moral Choice:} - The Romans are just following orders - Killing
them makes the party murderers in Rome's eyes - Can they find a way to
stop them without bloodshed?

\section{Scene 4: The Sacred Grove}\label{scene-4-the-sacred-grove}

Deep in the forest lies an ancient Germanic holy site---older than Rome,
older than memory.

\textbf{Read aloud:} \textgreater{} You emerge into a clearing unlike
any in the forest. Seven standing stones ring a granite altar, each
carved with spiraling runes that seem to shift in the firelight. The air
thrums with power---old power, wild power. The spear pulses in response,
eager and hungry. Thusnelda the druid kneels at the altar: ``Here was
where the God-Killer was forged. Here it must be unmade.''

\textbf{The Ritual Requirements:}

Thusnelda explains: 1. The spear must be placed on the altar 2. Blood of
a warrior must anoint it (1d6 damage, voluntary) 3. A sacrifice must be
made---something precious given willingly 4. The party must speak the
words: ``We unmake what was made in darkness''

\textbf{The Sacrifice:} - Can be a magic item - Can be a memory (lose
proficiency in a skill) - Can be an NPC ally who volunteers -
\textbf{Or} someone can take the corruption upon themselves (instant
corruption level 5)

\section{Scene 5: Mars Intervenes}\label{scene-5-mars-intervenes}

As the ritual begins, \textbf{Mars himself manifests}---not physically,
but as a divine presence.

\textbf{Encounter: Avatar of Mars} - Use \textbf{Fire Elemental} stat
block (CR 5) but reskin as divine wrath - Appears as a towering figure
of flame and fury - Speaks: ``THE SPEAR IS MINE! YOU DARE DESTROY WHAT I
CREATED?''

\textbf{Combat Options:} 1. \textbf{Fight:} Destroy the avatar (hard
battle) 2. \textbf{Complete the Ritual:} Hold it off while one person
finishes the ritual (skill challenge) 3. \textbf{Bargain:} Offer Mars
something else in exchange (what could a god want?)

\textbf{Skill Challenge: Complete the Ritual} - 4 successes before 3
failures - Each round, one person makes a check while others defend: -
Religion DC 15 (speak the words of unmaking) - Arcana DC 15 (channel the
grove's ancient magic) - Athletics DC 13 (hold the spear steady) -
Insight DC 14 (resist the spear's influence)

\section{Conclusion}\label{conclusion-1}

\textbf{If they destroy the spear:} - It shatters with a scream that
echoes across the spiritual realm - Mars' avatar dissipates, roaring:
``YOU HAVE MADE AN ENEMY OF WAR ITSELF!'' - The corruption fades from
anyone affected - Thusnelda: ``You have done what heroes of old could
not. But Mars does not forget.''

\textbf{If they fail or keep it:} - Mars claims the spear---or whoever
holds it - They must flee with a god's wrath upon them

\textbf{Advance to Level 6}

\section{DM Notes}\label{dm-notes-2}

\subsection{Guy Sclanders' ``Sacrifice
Principle''}\label{guy-sclanders-sacrifice-principle}

Victory should cost something meaningful. Players should feel they've
earned the win through sacrifice, not just dice rolls. The ritual's
sacrifice requirement embodies this perfectly.

\subsection{Branching Paths}\label{branching-paths}

This session has multiple valid outcomes based on player choices: - Path
A: Spear destroyed, Mars angry, party must deal with consequences - Path
B: Spear kept, corruption spreads, different finale ahead - Path C:
Bargain struck with Mars, new alliance formed

Track which path they take for Session 4's opening. \# Chapter 4: Return
to Rome

\section{Session Overview}\label{session-overview-3}

The party returns to Rome as either heroes or fugitives. They must
navigate the political intrigue of the Senate, discover who truly wanted
the spear, and prepare for the final confrontation.

\textbf{Level:} 6 → Advance to 7 at session end \textbf{Key Goal:}
Uncover the conspiracy and choose a side

\section{Opening: The Eternal City}\label{opening-the-eternal-city}

\textbf{If spear was destroyed:} \textgreater{} Rome rises before you
like a marble dream. The Capitoline Hill gleams white in the sun, the
Senate House stands proud, and the Colosseum roars with distant crowds.
But beneath the beauty, you sense rot. The gods are angry. Mars most of
all.

\textbf{If spear still exists:} \textgreater{} Rome beckons like a
serpent's smile. Somewhere in those marble halls, someone is searching
for the spear you carry. Someone powerful. Someone willing to kill for
it. The city that once felt like home now feels like a trap.

\section{Scene 1: Audience with the
Senate}\label{scene-1-audience-with-the-senate}

The party is summoned before a secret Senate committee in a private
chamber.

\textbf{Present:} - \textbf{Senator Gaius Cassius Brutus} (the primary
villain) - \textbf{Senator Marcus Tullius} (an ally) - \textbf{Vestal
Virgin Cornelia} (neutral observer) - \textbf{Tribune Lucius} (if alive)

\textbf{Senator Brutus' Proposal:}

He's smooth, charming, and utterly ruthless: - Congratulates them on
dealing with the ``frontier incident'' - Reveals he was behind the
Tribune's orders - Offers them wealth, land, citizenship for their
families - \textbf{In exchange:} They must retrieve the spear (if
destroyed) or give it to him (if it exists)

\textbf{His True Plan (DC 16 Insight to sense):} - He plans to use the
spear to assassinate the Emperor during the upcoming Triumph - Install
himself as dictator ``to restore the Republic'' - Believes Mars will
bless this if enough blood is spilled

\textbf{Senator Tullius' Warning:} - Pulls them aside afterward -
``Brutus is mad. He dreams of being Romulus reborn.'' - Offers to help
stop Brutus if they'll testify against him

\section{Scene 2: The City of
Shadows}\label{scene-2-the-city-of-shadows}

The party has 24 hours before the Emperor's Triumph. They can
investigate:

\textbf{Option A: Follow Brutus' Agents} - DC 14 Stealth to tail them
through Rome - Discover an underground temple to Mars - Cultists
preparing a ritual to resurrect the spear - Combat: 6 \textbf{Cult
Fanatics} (CR 2 each)

\textbf{Option B: Gather Evidence} - DC 15 Investigation in the Senate
archives - Find records of Brutus embezzling military funds - His
correspondences mention ``the God-Killer returns'' - Can use this to
turn other Senators against him

\textbf{Option C: Seek Divine Aid} - Visit the Temple of Jupiter - The
High Priest performs augury (costs 100 gp) - Jupiter speaks through
omens: ``The scales must balance. Blood for blood. Ambition for
humility.'' - Cryptic but helpful: suggests Mars can be appeased with a
worthy sacrifice

\textbf{Option D: Confront Brutus Directly} - Dramatic scene in his
villa - He monologues about Rome's corruption - ``The Republic died with
Caesar. I will resurrect it.'' - Can convince him to stand down (DC 18
Persuasion) or goad him into revealing plans

\section{Scene 3: The Underground
Temple}\label{scene-3-the-underground-temple}

Following clues leads to Brutus' secret temple beneath Rome.

\textbf{Read aloud:} \textgreater{} Beneath the city, deeper than the
sewers, you find it---a cavern carved into living rock. The walls run
red with ochre paint made to look like blood. An altar of black stone
dominates the center, and upon it\ldots{} a spear. Not the original---a
replica, but one infused with dark ritual magic. Around it, cultists
chant in Latin: ``Mars Ultor! Mars Ultor! Avenge us!''

\textbf{Encounter:} - \textbf{Senator Brutus} (use \textbf{Priest} stat
block but with +2 AC from magic armor) - 4 \textbf{Cult Fanatics} - 1
\textbf{Gladiator} (his personal bodyguard)

\textbf{The Ritual:}

If the party disrupted it in time: - Brutus rages: ``You've doomed Rome!
The gods demand blood!'' - He attacks with the replica spear (functions
as +1 magic spear, but not cursed)

If they're too late: - Brutus completes the ritual - Mars partially
possesses him (\textbf{Brutus transforms: HP doubled, +4 to hit, attacks
twice per turn}) - ``I AM THE HAND OF MARS! I WILL PURGE ROME IN FIRE
AND BLOOD!''

\section{Scene 4: The Emperor's
Triumph}\label{scene-4-the-emperors-triumph}

The climax occurs during the public Triumph---the Emperor's victory
parade.

\textbf{Setting:} The streets of Rome, thousands of citizens watching

\textbf{If Brutus escaped:} - He makes his move during the parade -
Attempts to kill the Emperor with the spear/replica - Party must stop
him publicly - Chase scene through crowds (skill challenge)

\textbf{If Brutus was captured:} - His remaining agents make the attempt
- Party must identify the assassin among the crowd - DC 16 Perception
each round (3 rounds before the Emperor is in danger)

\textbf{Public Combat:} - Fighting in front of Rome - Civilians scream
and flee (difficult terrain) - Praetorian Guards join the fight - Every
action is witnessed---reputation consequences

\section{Conclusion}\label{conclusion-2}

\textbf{Victory Outcomes:}

\textbf{If they saved the Emperor:} - Honored as heroes - Granted Roman
citizenship and land - But: Mars is still angry about the spear

\textbf{If they exposed the conspiracy:} - Senate is purged of
conspirators - They're celebrated but also feared - ``Those who topple
Senators can topple anyone''

\textbf{If Brutus escaped:} - He's now a fugitive with dangerous
knowledge - Sets up potential sequel - Mars' cult grows in the shadows

\textbf{Advance to Level 7}

\section{DM Notes}\label{dm-notes-3}

\subsection{Senator Brutus as Villain}\label{senator-brutus-as-villain}

Using Guy Sclanders' villain archetypes, Brutus is a ``Never Present''
villain until this session. The party has been fighting his plans, not
him directly. Now he's revealed---and he's not just evil, he's a
\textbf{tragic antagonist} who genuinely believes he's saving Rome.

\subsection{Moral Complexity}\label{moral-complexity}

Players might sympathize with Brutus' goals (ending corruption,
restoring the Republic) even if they disagree with his methods. This
creates dramatic tension---classic Sclanders storytelling.

\subsection{Setting Up the Finale}\label{setting-up-the-finale}

Whatever outcome, Session 5 will deal with Mars' wrath. The god of war
doesn't forgive, and the party has made an enemy of divinity itself. \#
Chapter 5: The Wrath of Mars

\section{Session Overview}\label{session-overview-4}

The final confrontation. Mars, god of war, demands satisfaction for the
destruction of his spear. The party must face divine judgment, make
ultimate sacrifices, and determine the fate of Rome itself.

\textbf{Level:} 7 (final level) \textbf{Key Goal:} Survive Mars'
judgment and end the campaign

\section{Opening: The Sky Bleeds Red}\label{opening-the-sky-bleeds-red}

Three days after the Triumph, omens appear across Rome.

\textbf{Read aloud:} \textgreater{} The sky turns red at noon. Not
sunset red---blood red. The Tiber runs backwards. Every dog in Rome
howls in unison. In the Forum, the statue of Mars begins to weep crimson
tears. The augurs declare what everyone already knows: the god of war is
coming, and he is angry. \textgreater{} \textgreater{} That night, you
all share the same dream. You stand in a vast arena of black sand
beneath a blood-red sky. Mars himself sits upon a throne of swords and
bones. His voice shakes your souls: ``YOU DESTROYED WHAT WAS MINE. NOW
PAY THE PRICE.''

\section{Scene 1: The Divine Summons}\label{scene-1-the-divine-summons}

The party awakens in a pocket dimension---Mars' personal arena, pulled
from the mortal world.

\textbf{The Arena:} - Circular, 200 feet across - Black sand that drinks
blood - Walls lined with the weapons of fallen heroes - No escape except
victory or death

\textbf{Mars' Judgment:}

The god appears---15 feet tall, clad in bronze armor, face hidden behind
a Corinthian helmet.

\textbf{His Voice (speak with gravitas):} \textgreater{} ``MORTALS. You
defied a god. You destroyed the Spear of Broken Oaths, forged in divine
fire and quenched in the blood of my son. For this crime, you will
either PROVE YOUR WORTH in combat, or DIE screaming. Choose your
champion---or face me as one.''

\textbf{The Challenge:} - Option A: One party member fights Mars in
single combat - Option B: Entire party fights Mars together (harder but
shared risk) - Option C: Attempt to bargain (what can you offer a god?)

\section{Scene 2: Trial by Combat}\label{scene-2-trial-by-combat}

\textbf{Option A: Champion vs.~Mars}

Mars summons a \textbf{Divine Champion} (use \textbf{Warlord} stat
block, CR 12 but adjust for single PC): - Actually a manifestation of
Mars' power, not the god himself - HP: 150, AC: 18 - Multiattack: Three
attacks with spear (+9 to hit, 1d8+5 piercing + 2d6 fire) -
\textbf{Divine Fury}: Once per turn, can make an extra attack -
\textbf{War God's Resilience}: Resistance to all damage except radiant

\textbf{Arena Mechanics:} - Weapons of fallen heroes line the walls (can
grab them as bonus action) - Each weapon grants different power: -
\textbf{Gladius of Achilles}: +2 weapon, advantage on attack rolls but
vulnerability to piercing - \textbf{Shield of Horatius}: +3 AC, can use
reaction to impose disadvantage on attacks - \textbf{Spear of Leonidas}:
Reach 10 ft, +1d10 damage if you haven't moved this turn

\textbf{Victory Condition:} - Reduce the champion to 0 HP - Or survive
10 rounds (Mars will respect endurance)

\textbf{Option B: Party vs.~Mars (Himself)}

\textbf{Mars, God of War} (Custom stat block): - HP: 300, AC: 20 -
Legendary Resistance (3/day) - \textbf{Multiattack}: Mars makes four
attacks - \textbf{God-Forged Blade}: +12 to hit, 2d8+7 slashing + 2d8
fire - \textbf{Divine Smite}: (Recharge 5-6) All creatures within 20
feet, DC 18 Con save, 8d6 fire damage (half on success) -
\textbf{Legendary Actions} (3/turn): - \textbf{Attack}: Make one weapon
attack - \textbf{War Cry}: Frightened condition, DC 16 Wisdom save -
\textbf{Summon Warriors}: 4 Shadows of the Fallen appear (use
\textbf{Specter} stats)

\textbf{Tactical Notes:} - Mars is arrogant; exploitable with clever
tactics - The arena weapons help but have costs - He respects
courage---reckless bravery can impress him

\section{Scene 3: The Bargain}\label{scene-3-the-bargain}

\textbf{If they attempt to bargain:}

Mars laughs: ``WHAT COULD MORTALS OFFER A GOD?''

\textbf{Possible Offers:} 1. \textbf{Service}: Become his champions,
fight in his name for a year and a day 2. \textbf{Sacrifice}: One party
member offers their life willingly 3. \textbf{Victory}: Win a war in
Mars' name (leads to future quest) 4. \textbf{Honor}: Promise to build a
new temple to Mars and spread his glory 5. \textbf{Knowledge}: Reveal
secrets of divine weapons (if they learned anything)

\textbf{Persuasion Check:} - DC 20 Persuasion or DC 18 Religion
(understanding what gods value) - Success: Mars accepts, modifies the
trial - Failure: ``YOU INSULT ME WITH EMPTY WORDS!'' (combat begins)

\textbf{Compelling Offers:} - ``Great Mars, we destroyed the spear not
to defy you, but because it was being misused by men unworthy of your
gift.'' - ``Let us prove our devotion by serving you---not as slaves,
but as honored champions.'' - ``The spear's corruption dishonored your
name. We purified your reputation.''

\section{Scene 4: The Price of
Victory}\label{scene-4-the-price-of-victory}

\textbf{After defeating/bargaining with Mars:}

The god is impressed (or satisfied):

\textbf{Mars' Decree:} \textgreater{} ``You have proven yourselves
worthy. Few mortals have stood before a god and lived. I grant you this:
your lives, and my grudging respect. But know this---you have interfered
with divine will. All actions have consequences.''

\textbf{The Cost (pick 2-3 based on campaign):}

\begin{enumerate}
\def\labelenumi{\arabic{enumi}.}
\item
  \textbf{Marked by War}: Each character bears Mars' brand (glowing red
  scar). They have advantage on Intimidation checks but disadvantage on
  Persuasion with peaceful folk. They are called to war whenever Mars
  wills it.
\item
  \textbf{The Burden}: One character (their choice) becomes Mars'
  Chosen. They gain a powerful boon (resistance to fire, extra attack,
  etc.) but must answer Mars' call once per year to fight in his name.
\item
  \textbf{The Debt}: They owe Mars a future favor---to be called in when
  he chooses. This could lead to a sequel campaign.
\item
  \textbf{The Transformation}: One willing character is transformed into
  a demigod servant of Mars---they leave the party to serve at his side
  (noble sacrifice ending for that PC).
\item
  \textbf{The Prophecy}: Mars reveals a dark prophecy about Rome's
  future: ``In three generations, Rome will burn. Only those marked by
  war can prevent it.'' (Sequel hook)
\end{enumerate}

\section{Scene 5: Return to Rome}\label{scene-5-return-to-rome}

The party awakens in the Temple of Mars in Rome, collapsed before the
altar.

\textbf{Epilogue Choices (based on their decisions throughout):}

\textbf{Heroic Ending:} - Hailed as saviors of Rome - Granted estates,
citizenship, military commissions - Senators seek their counsel - But
they carry Mars' mark---forever changed

\textbf{Pragmatic Ending:} - Quietly rewarded, given land on the
frontier - They leave Rome, knowing too much about its corruption -
Settle as veterans, with Mars' blessings and curses

\textbf{Tragic Ending:} - One or more died in Mars' arena - Survivors
honored but broken - Rome celebrates them as martyrs - Their sacrifice
is remembered in legend

\textbf{Rebellious Ending:} - They rejected both Mars and Rome - Now
hunted, they flee beyond the Empire - Freedom, but at the cost of home -
Legends grow of the ``God-Defiers''

\section{Final Scene: Where Are They
Now?}\label{final-scene-where-are-they-now}

Have each player describe where their character is one year later: -
What are they doing? - How has Mars' mark affected them? - What do they
remember most about this campaign? - What rumors do they hear about the
others?

\textbf{DM Final Narration:} \textgreater{} The Shadow of Mars is long.
What you did in those five sessions echoes across the Empire. Some call
you heroes. Some call you fools. Mars himself calls you\ldots{}
interesting. And in the divine realm, that is far more dangerous than
being forgotten. Your story is over. But legends never truly end.

\section{Campaign Conclusion}\label{campaign-conclusion}

Thank the players. Roll credits. Share highlights.

\textbf{Experience Wrap-Up:} - All characters end at level 7 - Award
bonus XP for exceptional roleplay - Distribute any remaining
treasure/rewards

\textbf{Sequel Hooks} (if desired): - Mars' prophecy about Rome's
burning - Brutus' cult surviving in shadows - Other gods taking interest
in the party - Germanic alliance or conflict - The spear's fragments
scattered across the world

\section{DM Notes: Guy Sclanders' Principles
Applied}\label{dm-notes-guy-sclanders-principles-applied}

\subsection{Master Plot}\label{master-plot}

The entire campaign followed the ``Spear's Journey'': Discovery →
Corruption → Destruction → Divine Judgment. Each session escalated the
stakes while maintaining player agency.

\subsection{Theme}\label{theme}

\textbf{Duty vs.~Free Will} - The party was caught between Rome's
orders, the gods' will, and their own conscience. Every major decision
forced them to choose what they truly valued.

\subsection{Villain Archetypes}\label{villain-archetypes}

\begin{itemize}
\tightlist
\item
  \textbf{Legate Corvinus}: Blunt Force Nemesis (direct opposition,
  Session 1-2)
\item
  \textbf{Senator Brutus}: Never Present (manipulating from shadows,
  revealed Session 4)
\item
  \textbf{Mars}: Divine Antagonist (ultimate challenge, Session 5)
\end{itemize}

\subsection{Player Agency}\label{player-agency}

Throughout, multiple paths were valid. No ``right'' answer---only
consequences. This is peak Sclanders design.

\subsection{Meaningful Sacrifice}\label{meaningful-sacrifice}

From the grove ritual to Mars' final price, victories cost something
precious. Players should feel they \emph{earned} their ending.

\begin{center}\rule{0.5\linewidth}{0.5pt}\end{center}

\section{Appendix: Quick Reference}\label{appendix-quick-reference}

\subsection{Key NPCs Quick Stats}\label{key-npcs-quick-stats}

\textbf{Legate Marcus Corvinus}: Use \textbf{Knight} (CR 3)
\textbf{Tribune Lucius}: Use \textbf{Noble} with +2 AC (CR 2)
\textbf{Centurion Varro}: Use \textbf{Veteran} (CR 3) \textbf{Augur
Cassia}: Use \textbf{Priest} (CR 2) \textbf{Vercingetorix}: Use
\textbf{Gladiator} (CR 5) \textbf{Thusnelda (Druid)}: Use \textbf{Druid}
(CR 2) \textbf{Senator Brutus}: Use \textbf{Priest} + magic armor (CR 3)

\subsection{Corruption Track
Reference}\label{corruption-track-reference}

\textbf{Level 0}: No corruption \textbf{Level 1}: Disturbing dreams
\textbf{Level 2}: Disadvantage on Wisdom saves \textbf{Level 3}: DC 12
Wisdom save to retreat \textbf{Level 4}: DC 14 Wisdom save to avoid
attacking allies \textbf{Level 5}: Fully corrupted, NPC

\subsection{Spear of Mars Properties}\label{spear-of-mars-properties}

\textbf{+2 weapon} \textbf{Damage}: 1d8 piercing + 1d8 necrotic
\textbf{Attunement Curse}: Disadvantage on Wisdom saves, increasing
aggression \textbf{Destruction}: Can only be destroyed at the sacred
grove via ritual

\begin{center}\rule{0.5\linewidth}{0.5pt}\end{center}

\textbf{Campaign Complete! May Mars smile upon your future games!}


\backmatter


\end{document}
